\chapter{Conclusion}
\section{Summary}
This first stage of the project has successfully established a stable foundation, covering requirement analysis, system design, technology selection, and core implementation. The system now demonstrates a functional workflow and provides a solid base for expansion in the next stages. Although the initial version focuses primarily on delivering essential features, it also opens the door for a wide range of improvements that will significantly enhance usability, performance, and the overall learning experience.

\section{Future Work}
Moving forward to the second stage, the project will focus on severa major developments:
\begin{itemize}
    \item \textbf{Migration of storage services to AWS S3:} This transition will improve reliability, scalability, and integration capabilities, ensuring the platform can support increasing amounts of user-generated content.
    \item \textbf{Deployment to a production-ready environment:} In the next phase, the system will be deployed on a cloud platform, enabling real-world accessibility, automated scaling, and continuous integration/continuous delivery (CI/CD). 
    \item Most importantly, we will introduce several new functions to support learner better beside our main "Trascript and Summarize" function:
    \begin{itemize}
        \item Automatic mind-map generation from transcripts or uploaded documents.
        \item Integrated AI chatbot able to answer questions, summarize content, or guide learners using the processed data.
        \item Flashcardto reinforce learning based on extracted topics.
    \end{itemize}
\end{itemize}
Together, these future improvements will transform the current prototype into a comprehensive learning-support system capable of not only processing content but also generating meaningful learning materials. With a clear roadmap and a strong technical foundation, \textit{Quizum} is positioned to evolve into a highly scalable, intelligent, and user-friendly platform in the next development stages.

